% 核心观点：使正文中的通用设计可以追溯到标签、任务实例和执行设置。
% 叙述逻辑：现有数据来源及缺失证据 -> 待实现审核协议 -> 已有关系实例 -> 已有调度与槽位映射。
% 本附录不将新增审核、校准或扰动实验写成已完成模块或结果。
\section{Method Details and Implementation Status}
\label{app:methoddetails}
\subsection{数据来源与划分口径}
\label{app:datascope}
本附录以本地保存的Qtoken-FS及其SF-C-GRU配置为可核对实例，
不将其自动认定为最终实验版本。
基础训练配置指向bettersetup-v5，入口读取传入数据根目录的全部轨迹；
配置没有列出训练、验证和测试episode ID，上游是否独立划分仍需原始清单确认。
front与side分割清单分别指向bettersetup-v4和bettersetup-v5，
各记载15118个训练样本、2819个验证样本，但两个列表均仅保留前2000项。
其中的视频文件索引也不能直接等同于episode ID。
这些记录能证明存在分割器内部划分，不能证明它与策略使用同一划分。

残差缓存生成器按任务分层，以固定种子划分约80\%训练和20\%验证轨迹，
并明确将验证范围限定为残差头；冻结基础策略可能已见全部轨迹。
本地保存了SF-C-GRU的缓存指纹，但未找到该缓存的原始划分文件，
故不能重建其确切episode名单，也不能由种子猜测名单。
各阶段证据及缺失项登记在随稿的\texttt{data\_split\_audit.json}；
其未知名单使用空值，不能当作已经完成的统一划分。

后续全流程独立实验须先建立共同的轨迹清单，记录原始采集ID、数据版本、
各阶段用途及派生文件映射，再由提示器、分割器、基础策略和残差训练复用。
若固定已有上游模型，只重新划分残差数据，则结果继续限定为残差头独立验证。
当前证据既不足以断言最终测试泄漏，也不足以声称全流程未见。

\subsection{分割质量的实现边界}
\label{app:seguncertainty}
\label{app:visibilityextension}
当前流程使用掩码质量诊断及关系监督加权，未实施概率温度校准、扰动集成或人工三态可见性标签接入。
面积项仍将空掩码作为零面积监督，故其可靠性必须通过独立人工审核评价。
若后续区分确认可见、确认不可见与未确定，须把审核状态实际接入关系有效标记和损失权重，
而非仅增加一个诊断分数。改变软语义概率的校准还需同步重训或验证基础策略。
这些扩展不是当前方法组成部分，也不用于支持已有结果。

\subsection{关系查询的任务实例}
已有查询配置采用三组front图像关系：目标可见面积占比、目标到目标区域的质心距离，
以及目标质心相对工具端点的位置。
在宽$W$、高$H_I$的front图像中，
离线硬标签的目标、区域和工具掩码分别为$M_O,M_R,M_T$，
非空掩码的像素质心为$c_O,c_R,c_T$。
三个关系的具体算子为
\begin{equation}
 \begin{aligned}
 \widetilde G^{(1)}&=|M_O|/(WH_I),\\
 \widetilde G^{(2)}&=\norm{c_O-c_R}/
       \sqrt{(W-1)^2+(H_I-1)^2},\\
 \widetilde G^{(3)}&=\clip(\bar c_O-\bar e_{T,L},-1,1).
 \end{aligned}
 \label{eq:methodgeometry}
\end{equation}
其中，$\bar c=(c_x/(W-1),c_y/(H_I-1))$为归一化质心，
$\bar e_{T,L}$为归一化工具端点，$\clip$逐分量裁剪。
坐标原点位于图像左上角，横轴向右、纵轴向下；
因此$\widetilde G^{(3)}$的两个分量分别为目标相对端点向右、向下的偏移。

工具端点由归一化工具像素的协方差矩阵主特征向量$a_T$确定。
对工具像素$\bar\xi$，计算投影$a_T^\top(\bar\xi-\bar c_T)$的极值$\alpha_-,\alpha_+$，
得到$\bar e_-=\bar c_T+\alpha_-a_T$与$\bar e_+=\bar c_T+\alpha_+a_T$。
取横坐标较小者为$\bar e_{T,L}$，横坐标相同时沿用实现选择$\bar e_-$。
它不自动等同于物理接触端；方向变化可能改变这一约定的适用性，
须在任务实例中核验。
区域质心距离也不等同于物体被完整包含，因此任务完成评价仍采用独立定义。

当前实现根据掩码质量和几何有效性确定监督权重。
面积组的有效标记恒为1，空目标掩码直接生成零面积标签；
位置组要求对应目标和区域或工具掩码的像素质量均大于$10^{-6}$。
实现以占位数值填充无效关系，但由零权重排除其损失。
它尚未读取附录~\ref{app:visibilityextension}中的人工审核状态。
当前三组输出维度为1、1、2，每个查询为512维，
采用$\beta_G=0.01$的Smooth L1损失。
已有实现将关系所涉及类别的质量分数最小值记为$\kappa_{nj}$，
以有效标记$m_{nj}^{G}$构造组内权重
$w_{njc}=m_{nj}^{G}\min(\kappa_{nj}/0.95,1)$。
这里$\kappa_{nj}$为启发式质量分数，不是校准后的正确概率。
分组归一化主要改变样本的相对贡献；
当分母不受截断影响时，全组同倍降权会被分母抵消。

Qtoken-FS保存的类别颜色见表~\ref{tab:semanticpalette}。
实现中的palette为浮点数；表内8-bit RGB各通道除以255得到式\eqref{eq:methodsoftsem}的
$e_c^{\mathrm{sem}}\in[0,1]^3$，再按类别概率加权，不先取argmax或量化为整数。
\begin{table}[t]
\caption{Qtoken-FS记录中的语义颜色与类别顺序（8-bit RGB），不是最终配置的默认值。}
\label{tab:semanticpalette}
\centering
\begin{tabular}{lll}
\toprule
类别 & RGB编码 & 视角\\
\midrule
background & $(0,0,0)$ & front/side\\
occluder & $(64,160,255)$ & front/side\\
object & $(255,105,97)$ & front/side\\
region & $(119,221,119)$ & front/side\\
tool & $(255,209,102)$ & front/side\\
leftarm & $(234,146,199)$ & front\\
rightarm & $(173,214,101)$ & front\\
\bottomrule
\end{tabular}
\end{table}
通道按表~\ref{tab:semanticpalette}中各视角可用类别的顺序排列。
该实例的front含机械臂类别，side不含；不能笼统称为所有视角均未分割机械臂。
保存的语义温度为1，未拟合附录中的概率校准参数。
已有基础网络使用完整视野的$480\times270$图像、
共享ResNet18骨干及ACT投影与位置编码；
双视角含四路视觉输入，单视角含两路。
当前该配置不额外加入视角或模态embedding，也不包含阶段分类或停止头。
保存配置中的动作损失系数和关系损失系数均为1；
ACT默认KL系数为10，训练日志亦满足
$L_{\mathrm{action}}=L_{\mathrm{L1}}+10L_{\mathrm{KL}}$。
因此，该记录对应式\eqref{eq:methodbaseloss}中的$\lambda_{\mathrm{KL}}=10$、
$\lambda_{\mathrm{rel}}=1$，而不是待猜测的默认值。
冻结分割器在基础策略阶段的分割损失为零。
这些值对应现有Qtoken-FS记录；最终采用不同checkpoint时须重新核对。

\subsubsection{ACT潜变量的训练与部署}
基础训练启用CVAE时，由本体观测与示范动作序列构造后验并采样潜变量；
其投影token与关系查询、图像和本体token进入同一编码器。
关系辅助头读取该编码器的查询输出，故其训练损失也受后验潜变量条件影响。
测试时潜变量置零，沿用ACT的约定\cite{zhao2023act}。
残差缓存适配器将冻结基础策略置于\texttt{eval()}模式，
并仅向动作块预测接口提供图像与本体输入，缓存查询也来自这次部署模式前向计算。
后续关系误差评价须使用相同模式，另报告训练后验条件与零潜变量条件的差异；
不将训练关系损失直接作为部署关系精度，也不在测量前断言存在动作信息捷径。

\subsection{残差训练与执行设置}
已有完整反馈变体使用双视角RGB与缓存关系查询。
新图像完整缩放至$320\times180$，经冻结骨干编码，
每视角特征池化为$2\times3\times512$；
仅RGB与额外分割描述作为不同输入变体。
当前记录使用128维投影、单层hidden128的GRU和8条历史记录。
目标反馈频率为5Hz，等间隔满窗口的跨度约1.4秒；
在线间隔受共享线程调度影响，须由实际时间戳报告。每次从零状态重算窗口，
而不是跨episode保持无限递归记忆。
当前信息MLP使用同一记录格式中的最新一条，保留计划、本体和即时运动上下文。

监督指令来自leader，观测关节来自follower。
动作与残差使用数据集原生命令单位，关节顺序、动作维度及单位须与部署一致。
每维训练集动作标准差以$10^{-3}$为下限作为$\sigma_c$；
残差限幅取训练集基础动作误差绝对值的95分位数，再裁剪到$[0.05,5]$原生单位。
这些值是当前实现约定，不是跨平台通用安全阈值。
残差训练采用$\beta_r=0.1$、$\lambda_r=10^{-3}$，
AdamW学习率$10^{-4}$、batch128、最多50轮，
验证8轮不改善早停，以验证集尺度归一化MAE选择checkpoint。

SF-C-GRU保存配置的控制频率为30Hz，基础动作块为60步，基础更新目标间隔为30步；
反馈目标间隔为6步，预留3步延迟，预测9步修正。
在第6帧获取的观察对应第9--17帧，第12帧观察对应第15--23帧，
重叠3帧，使用首槽的端到端预算为100ms。
该记录中的观测年龄、残差年龄、历史断流阈值及视角时间差上限分别为
0.15秒、0.4秒、0.3秒和0.04秒。
历史断流阈值允许正常的0.2秒目标间隔，若实际资源等待超过0.3秒，历史仍会被清空。
离线缓存按固定步长及预设延迟构造样本，
没有重演在线基础任务优先与共享线程忙碌造成的反馈等待。
因此，5Hz与1.4秒仅为目标设置及其名义窗口跨度，不是实测执行频率。

完整时效分析应记录每个候选的观察时间、提交与完成时间、计划编号、目标槽位、
接收状态及最终执行状态；现有日志字段不能覆盖的量须另补采集，不回填推测值。
延迟统计包含采集、预处理、编码、排队、GRU和指令发布，
不以网络计算耗时代替端到端时效。
无残差时复用有效基础计划；计划过期时当前桥接代码抛出异常并退出策略。
异常之后的硬件处理仍需单独验证，不将抛出异常视为安全保证。

\subsubsection{时间戳与目标槽位}
\label{app:slottiming}
采用控制频率$f_{\mathrm{ctrl}}$和运行起点$s_{\mathrm{origin}}$，
第$k$次反馈的图像参考时间为
$a_k=\min_v(s_k^v-s_{\mathrm{origin}})$，
其中$s_k^v$为视角$v$的采集完成时间戳。
预留延迟$d_{\mathrm{fb}}$以控制步计，候选槽位按当前实现映射为
\begin{equation}
 t_{k,j}=\operatorname{round}\!\left[
 f_{\mathrm{ctrl}}(a_k+d_{\mathrm{fb}}/f_{\mathrm{ctrl}})\right]+j.
 \label{eq:methodslots}
\end{equation}
其中$j=0,\ldots,H_r-1$，$\operatorname{round}$采用代码中的Python最近整数规则。
在与控制网格对齐的离线帧上，$n_k=f_{\mathrm{ctrl}}a_k$为整数，
该式化为$t_{k,j}=n_k+d_{\mathrm{fb}}+j$。
在线时间戳不一定落在网格上，应保留式\eqref{eq:methodslots}的舍入顺序。

执行循环以相对墙钟$s$确定全局槽位
$\operatorname{round}(f_{\mathrm{ctrl}}s)$；
基础动作块以自己的观察时间$a_p$为起点，读取相对索引
$\operatorname{round}(f_{\mathrm{ctrl}}(s-a_p))$。
这两次舍入与候选槽位映射分别记录，不直接相减已舍入时间戳替代相对索引。
实际控制循环可能跳过槽位，未执行槽位的候选不顺延使用。
基础预测与反馈共用单工作线程和一个待完成任务，
到期基础更新优先；反馈在空闲且观测有效时才提交。

\subsection{实现边界与复现记录}
已有代码提供掩码时序诊断、软颜色语义输入、
front关系查询监督、固定窗口GRU/MLP残差及计划槽位校验。
新增的人工校准划分、温度拟合、扰动关系诊断与其完整审核报告尚未实现，
原有模型也不能视为已经使用这些机制。
最终实验须记录分割与策略权重、颜色及类别表、数据划分、
关系定义、残差输入变体和全部执行配置，
并区分已有结果与采用新质量流程重新训练、重新验证的结果。
